\input{../tex-inputs/latex-leseansicht-vorspann}

\section[Gustav Schwarzkopf an Arthur Schnitzler, {[}19. 6. 1903{]}]{L04310 Gustav Schwarzkopf an Arthur Schnitzler, {[}19. 6. 1903{]}}
\nopagebreak\mylabel{L04310v}
\rehead{ }\normalsize\beginnumbering\briefempfaengerindex{Schnitzler, Arthur@\textsc{Schnitzler, Arthur}!zzzSchwarzkopf, Gustav@\emph{von Gustav Schwarzkopf}!1903-06-192@{{[}19. 6. 1903{]}}|(be}
\toendnotes[C]{\smallbreak\pagebreak[2]}
\correspDesc{Versand  durch Gustav Schwarzkopf am [19. 6. 1903] in Wien
\newline{}Erhalt  durch Arthur Schnitzler am [19. 6. 1903?] in Wien}\toendnotes[C]{\smallbreak}
\Standort{CUL, Schnitzler, B 96a.}
\physDesc{Brief, 1 Blatt, 2 Seiten, 226 Zeichen
\newline{}Handschrift: schwarze Tinte, deutsche Kurrent
\newline{}Schnitzler: mit Bleistift datiert: »19/6 903.« }\toendnotes[C]{\smallbreak}
\pstart
           \noindent{}{\pb}Lieber Arthur, beſten
               Dank, aber es iſt mir leider nicht möglich. Ich habe zugeſagt heute zu einer
               \label{K_L04310-1v}\edtext{Abſchiedsfeier zu ko{\geminationm}en, die Herr \textsc{Eppens\pwindex{Eppens, Otto 16.\,12.\,1860 Basel – 1921 Hamburg@\textsc{Eppens, Otto} (16.\,12.\,1860 Basel – 1921 Hamburg), \emph{Schauspieler}|pw}}}{\lemma{\textnormal{\emph{Abschiedsfeier … Eppens}}}\Cendnote{\textnormal{Nachdem er 13 Jahre am \emph{Deutschen Volkstheater}\orgindex{Volkstheater@Volkstheater|pwk} engagiert gewesen
               war, gab er am 19. 6. 1901 seinen letzten Auftritt auf dieser Bühne und übersiedelte
                  nach Hamburg\oindex{Hamburg@\textbf{Hamburg}|pwk}.}}}\label{K_L04310-1} gibt und beim »ſilbernen Brunnen\oindex{Wien@\textbf{Wien}!IX., Alsergrund@\textbf{IX., Alsergrund}!Zum silbernen Brunnen@\textbf{Zum silbernen Brunnen}, \emph{Gastgewerbegebäude}|pw}« abgehalten {\pb}wird. Sie und Frl Olga\pwindex{Schnitzler, Olga 17.\,1.\,1882 Wien – 13.\,1.\,1970 Lugano@\textsc{Schnitzler, Olga} (17.\,1.\,1882 Wien – 13.\,1.\,1970 Lugano), \emph{Schauspielerin, Sängerin}|pw} herzlich grüßend\pend
           
\pstart
           Ihr{\\[\baselineskip]}\spacefill\mbox{Gustav}\pend
           \leftskip=0em{}\selectlanguage{ngerman}\endnumbering\briefempfaengerindex{Schnitzler, Arthur@\textsc{Schnitzler, Arthur}!zzzSchwarzkopf, Gustav@\emph{von Gustav Schwarzkopf}!1903-06-192@{{[}19. 6. 1903{]}}|)be}\mylabel{L04310h}
\begin{anhang}
\end{anhang}\newcommand{\dateiname}{L04310}\newcommand{\titel}{Gustav Schwarzkopf an Arthur Schnitzler, [19. 6. 1903]}\newcommand{\editorInnen}{Herausgegeben von Jahnke, SelmaMüller, Martin Anton}\input{../tex-inputs/latex-leseansicht-abspann}
